\documentclass[10pt,a4paper]{amsart}
\usepackage[margin=19mm]{geometry}
\usepackage{amsmath,amssymb,mathtools}
\usepackage[scr=rsfs,cal=euler]{mathalfa}
\usepackage{xfrac}
\usepackage[unicode,hidelinks,bookmarks=false]{hyperref}
\newcommand{\ie}{i.e.\ }
\newcommand{\cf}{cf.\ }
\newcommand{\resp}{respectively\ }
\newcommand{\Subsection}{\subsection}
\providecommand{\nobreakdash}{\nobreak\hbox{-}}
\setlength{\emergencystretch}{2em}
\multlinegap0pt
\usepackage{xcolor}
\newtheorem{thm}{Theorem}
\def\Sp{\Sigma_\mathrm{pred}}
\def\Sph{\hat{\Sigma}_\mathrm{pred}}
\def\dy{d_\mathrm{y}}
\def\mph{\hat{\mu}_\mathrm{pred}}
\def\ufr{\mathrm{u}_\mathrm{f}}
\def\uft{\tilde{u}_\mathrm{f}}
\def\yr{y_\mathrm{ref}}
\title{\LARGE \bf Gaussian behaviors: representations and data-driven control}
\author{Andr\'as Sasfi, Ivan Markovsky, Alberto Padoan, Florian D\"orfler}
\date{Source: arXiv:2504.15838v2 (CC BY 4.0)}
\begin{document}
\maketitle

\noindent\emph{Source and licence.} \LARGE \bf Gaussian behaviors: representations and data-driven control. Version arXiv:2504.15838v2 (CC BY 4.0), distributed by arXiv under the Creative Commons Attribution 4.0 International licence, http://creativecommons.org/licenses/by/4.0/, as recorded in the arXiv licence metadata for this exact version and shown on the abstract page https://arxiv.org/abs/2504.15838v2. Full licence notice: \texttt{licenses/arxiv-2504.15838-CC-BY-4.0.txt}.\par\smallskip

\noindent\textit{Editorial scope.} Counted unit: the article's thm thm:final (source0.tex lines 571--582). The article's complete original proof of it is delivered with it (source0.tex lines 583--586, one \texttt{proof} environment of 4 lines). No mathematics is written, repaired or re-derived: the statement and the proof are the article's own lines, in the article's own notation, and no retained line is substituted or re-flowed.  Retained uncounted as the source-local context the proof needs: lines 562--567.  Nothing was omitted from any retained span, so the package carries no omission record.  No retained line contains a citation, so the wrapper carries no bibliography and no bibliography entry.  No retained line contains a figure environment, an image-inclusion command or drawing code, so no image is delivered and the package declares no asset dependency.  Editorial formatting only, in addition: an AMS \texttt{amsart} preamble that restates the article's own declarations and macros used by the retained lines, namely the environments \texttt{thm} on separate counters, together with the standard packages those lines need.  The retained environments print in this document as Theorem 1 in order of appearance, whereas the article prints the same items with the numbering of its own full document.  The complete native source of the article as distributed in the arXiv e-print for this exact version is bundled unchanged as \texttt{sources/176900/source0.tex}.
\begin{align} \label{eq:min-max-constrained}
    \begin{split}
    & \min_{\ufr \in \mathcal{U}} \max_{\mu\in\mathbb{R}^{\dy}} \;  \mathbb{E}_{\mathcal{N}(\mu,\Sp)}[\tilde{J}(\ufr)] \\
    & \quad  \mathrm{s.t.} \; D_{KL}\left(\mathcal{N}(\mu,\Sp)~\big\|~\mathcal{N}\left(\mph,\Sph\right) \right) \leq \epsilon.
    \end{split}
\end{align}

\begin{thm} \label{thm:final}
    The cost of~\eqref{eq:min-max-constrained} is upper bounded by
    \begin{align} \label{eq:final_opt}
        & \mathbb{E}_{\mathcal{N}(\mu^\star,\Sp)}[\tilde{J}(\uft)] \\
    & -\lambda\left( D_{KL}\left(\mathcal{N}(\mu^\star,\Sp)~\big\|~\mathcal{N}(\mph,\Sph)\right) - \epsilon\right), \nonumber
    \end{align}
    where
    \begin{align*} 
        \mu^\star := (\lambda \Sph^{-1}-Q)^{-1}(\lambda\Sph^{-1}\mph -Q\yr),
    \end{align*}
    for any $\uft$ and any $\lambda\geq\lambda_0$, where $\lambda_0>0$ is such that $\lambda_0\Sph^{-1}-Q\in\mathbb{S}_{>0}$.
\end{thm}

\begin{proof}
    The inner problem in~\eqref{eq:min-max-constrained} can be expressed as $\min_{\mu}-\mathbb{E}_{\mathcal{N}(\mu,\Sp)}[\tilde{J}(\ufr)]$ subject to the constraint on $D_{KL}$, which is lower bounded by its dual function. Thus, the negative of the dual, given in~\eqref{eq:final_opt}, is an upper bound on the inner problem in~\eqref{eq:min-max-constrained}.
    Since~\eqref{eq:min-max-constrained} is minimized over $\ufr$, its cost is lower than~\eqref{eq:final_opt} for any $\uft$.
\end{proof}

\end{document}
